\chapter{Introduction}
\label{chapter:introduction}

Robots are expected to navigate multiple and diverse environments. To achieve this, they need situational awareness and the autonomy to make proper decisions. Computer vision aims to enable machines to interpret the world that surrounds them. By processing images and videos, robots can extract information to recognize object geometries, textures, materials, etc. Although visual analysis allows a machine to build internal models containing vast knowledge of its environment, this does not mean that its perception is complete. Robotic perception gives a machine the ability to understand its surroundings to the degree that it can safely interact with the world \cite{RobotV}. Such understanding requires other techniques in addition to computer vision to obtain a complete sensory experience of the environment. Acoustic sensing and analysis can enhance machine perception by following an approach similar to that seen in nature with animals that use binaural acoustic perception: the environment can be analyzed to detect the source of a particular sound and focus visual attention on that region. However, this capability is not only an extension of vision. In some conditions, animals rely entirely on acoustic information to operate, as in the echolocation employed by bats, birds, dolphins, and whales, which hear the echoes of acoustic chirps to gain spatial perception of their surroundings and navigate complex layouts \cite{echoloco}. 

 Traditional computer vision techniques that employ sensing devices such as cameras and lidars provide high-fidelity information to navigation systems, but they become unreliable in scenarios with poor lighting or the presence of dust, smoke, and fog \cite{cat1}. Another aspect to consider is the cost of the sensing device. Acoustic perception offers an advantage in this regard due to the low cost of commercial electret or MEMS microphones compared to optical cameras and lidars, which makes it desirable as an alternative to traditional sensors. Recent developments in machine learning, such as those shown in \cite{bat1, cat1, echosense}, have demonstrated the great potential of neural networks to infer spatial and geometrical information about surrounding objects by analyzing acoustic echoes in the audible spectrum. The manufacturing sector, under the Industry 4.0 transformation, is increasingly adopting automation and smart monitoring technologies, which makes it an important field of application for these new methods. This is especially true for the supervision of machinery during operation, particularly machines on production floors that use robotic arms and computer numerical control actuators to manipulate objects in confined spaces \cite{mordor}. In such conditions, acoustic sensing is a good approach to gain insight into the overall status of the system, because the Room Impulse Response acts as a transfer function that encodes information about the environment setup. The acoustic response of a room can be used together with artificial neural networks, such as convolutional neural networks or recurrent neural networks, to learn from this response and recognize characteristics and properties of the room echoes, and thereby recognize spaces and objects within it. Thus, acoustic perception together with neural networks can improve the perception of machinery operating in complex spaces.  




%\newpage

\section{Problem}

The use of acoustic signals to approximate the size and shape of objects has been considerably improved by the adoption of machine learning techniques. An important part of the work in this field was developed for indoor navigation systems, where the objective is to enhance the perception of robots by using acoustic analysis to build low-resolution depth images of an unknown environment. This work presents a different point of view. First, the environment is completely known a priori, but an object changes its position inside it, and its location must be estimated with a certain accuracy each time. Second, in contrast to the large rooms in which those systems are intended to navigate, the object location must be estimated at short ranges, where multi-path echoes can be received at the same time. This creates a superposition that is difficult to resolve and can cause those systems to malfunction \cite{bat1}. To our knowledge, no existing work estimates the position of an object through acoustic analysis under the circumstances described. To achieve acoustic spatial estimation in this scenario, we propose two key changes to the processing pipeline. The first is in the beamforming stage, to increase the spatial discrimination of the origin of the echoes. The second is a change in the type of neural network used to infer the location, employing a network suited to directly process sequences of acoustic features instead of the images required by Convolutional Neural Networks. Our purpose is to demonstrate that it is possible to locate an object in a small room (e.g., 1 m side length) using acoustic signal processing with neural networks.

\section{Document Structure}

This document present in \textbf{Chapter 1} a brief introduction to the topic and is stated the problem we intent to address with the development of this thesis. \textbf{Chapter 2} contains the theoretical background and state of the art in the field. \textbf{Chapter 3} details the hypothesis, objectives and scope of the work. \textbf{Chapter 4} shows the research proposal and methodology. \textbf{Chapter 5} present the results obtained from the experiments conducted. \textbf{Chapter 6} discusses the findings and their implications. \textbf{Chapter 7} concludes the thesis and outlines future work.






